\section{Canonical marked widths and the strict unbuffered defect}
\label{sec:wave-proofs}
The word \emph{canonical} here describes the
specified base/crest \emph{unbuffered A.1 Wave class}.
It does not identify KKL's distinct
buffered canonical WAD of Appendix A.2.3.

\begin{definition}[Extremal width and three independent suprema]
For a family $\Gamma$ of locally rectifiable proper
curves in a marked Riemann surface $Y$, set
\begin{equation}\label{eq:width}
 \W_Y(\Gamma)=
 \inf\left\{\int_Y\rho^2\,dA:
    \rho\ge0,\ \int_\gamma\rho\,|dz|\ge1
       \text{ for every }\gamma\in\Gamma\right\}.
\end{equation}
If $Y=(0,L)\times(0,H)$ and $\Gamma$ consists
of the vertical leaves, $\W_Y(\Gamma)=L/H$.
Put $P_p=\sup\W(\mathcal S)$ over \emph{all}
KKL A.1 Waves on $\Om_p=U_p\setminus\T_p$
with base $(I_9,I_0,I_1)$ and crest
containing $I_0$.  Independently put
$C_p^+$ and $C_p^-$ equal to the suprema
of all A.1 Waves on
$\Om'_p=V_p\setminus\T'_p$
with bases $(S_1,S_2,S_3)$ and
$(S_{17},S_{18},S_{19})$, and crests
$S_2$ and $S_{18}$, respectively.
Reflection across $\R$ identifies
the two child classes isometrically,
hence $C_p^+=C_p^-=C_p$.
No attainment, selection of a
common parent, or bijection of
all arbitrary parent and child
families is presumed.
\end{definition}

\begin{lemma}[Physical marked-closure bridge]\label{lem:M}
The unbuffered marked surfaces used in the
canonical-sandwich and strict-gain calculations
are \emph{identically} $\Om_p$ and $\Om'_p$ as
embedded complex plane domains.  By
Theorems~\ref{thm:marking} and \ref{thm:incidence},
the identity extends to their physically
marked KKL Carath\'eodory ideal closures.
It therefore bijects \emph{every} A.1
proper-path class and Wave, with base,
ends and crest, and preserves its conformal
extremal width.  Thus the earlier
family-specific identification hypothesis
$M$ is discharged \emph{on these two
selected physical slits}.  No analogous
claim is made for a different physical
puzzle surface or for the four-hole WAD.
\end{lemma}
\begin{proof}
The point sets and complex coordinates
are exactly $U_p\setminus[x_{2p},x_p]$
and $V_p\setminus[-x_p,x_p]$, not
homeomorphic substitutes.  The
Carath\'eodory identification in
Theorem~\ref{thm:marking}
matches every marked side and every
outer whole-component circle.
KKL Appendix A.1 defines properness,
the forbidden length-at-most-two
subordination, relative-endpoint
homotopy into \emph{every} ideal collar,
simple laminations, their smallest
base multi-interval and its crest.
All these conditions are invariant
under the \emph{same marked-closure
identity} in both directions.
For a conformal coordinate map
$w=\phi(z)$ an admissible density
transforms as
$\rho_w(\phi(z))=\rho_z(z)/|\phi'(z)|$;
the curve lengths and area integrals
in \eqref{eq:width} are unchanged.
The identification here is the identity
and the auxiliary positive rescaling
$\xi=-A_pz$ is conformal.
Thus both classes and their
\emph{actual} widths are equal.
A bare topological equivalence
would not give the width equality.
\end{proof}

\begin{theorem}[Canonical marked Wave sandwich]
\label{thm:sandwich}
For every $p\ge3$ put
\[
 m_p=\frac{x_{2p}+x_{5p}}2,\quad
 r_{0,p}=\frac{x_{5p}-x_{2p}}2,\quad
 r_{1,p}=\min\{-\alpha_p-m_p,\ x_{4p}+m_p\},\quad
 W_p^{\rm rad}=\frac1{2\pi}\log\frac{r_{1,p}}{r_{0,p}}.
\]
Write $P_{p,2}$ for the smaller parent
supremum requiring that \emph{every}
leaf has the two named surviving proper
A.3.2 children.  For the full lifted-path
envelopes $Q_p,Q(L_{+,p}),Q_{\rm plane,p}$
defined in the proof,
\begin{equation}\label{eq:sandwich}
 W_p^{\rm rad}\le P_{p,2}\le P_p
 \le Q_p\le Q(L_{+,p})\le Q_{\rm plane,p}.
\end{equation}
For $p\ge13$,
\[
 0.26<P_p<0.78,\qquad
 0.26<\liminf_{p\to\infty}P_p
 \le\limsup_{p\to\infty}P_p<0.78.
\]
The limiting inequalities \emph{do not}
assert convergence or equality to an
extremal quadrilateral.
\end{theorem}
\begin{proof}
By \eqref{eq:order1}--\eqref{eq:order2},
$0<r_0<r_1$ (we suppress $p$).
For $r_0<s<r_1$ consider the open
circle $\gamma_s(t)=m+s e^{2\pi it}$,
$0<t<1$.  Its real right meeting
$q_s=m+s$ satisfies
$x_{5p}<q_s<-\alpha_p<x_{3p}$,
while its left meeting $w_s=m-s$
lies in $(-x_{4p},x_{2p})$.
Thus the \emph{only} removed target
slit point is the right meeting,
whose two accesses are precisely
$I_9,I_1$.  The curve goes around
the \emph{left} tip, so its
relative-endpoint class is
subordinate to the three-atom
arc $(I_9,I_0,I_1)$ and cannot
be subordinate to any one-
or two-atom arc.  This checks
KKL A.1 properness and the crest,
rather than assuming every
circle gives a Wave.

The complete closed radial gate disk
lies in the \emph{central} $V_p$.
For $3\le p\le12$ the whole-disk
bound and real connector are
certified by the finite Rouch\'e
interval checker; it is not a
test of one radius.  For $p\ge13$,
the normalized $h_p(w)=w^2+\tau_p+e_p(w)$
bounds the entire gate by
$|A_pz|<2$, and
\[
 |G_p(z)|<\frac7{|A_p|}
 <\frac{140}{3}4^{-p}
 <\frac35\,2^{-p}<R_p.
\]
Since the gate center belongs to
the certified central real pullback,
connectedness of the gate disk
places the whole disk in $V_p$.
The same circle in $\Om'_p$
is cut by the \emph{one} newly
attached real interval into two
proper open semicircles.  Their
source endpoints are $(S_3,S_1)$
and $(S_{19},S_{17})$, so both
survive the strict $>2$ atom
filter and meet distinct
immersion blocks.

The log chart
$z\mapsto\log((z-m_p)/r_{0,p})$
identifies the cut collar with
$(0,\log(r_1/r_0))\times(0,2\pi)$.
The vertical-leaf extremal density
and Cauchy--Schwarz give width
$\log(r_1/r_0)/(2\pi)$.
The two semicircle children have
angular height $\pi$, so each has
\emph{exactly} twice that width.
This is a width identity for
\emph{these chosen leaves only}.
The A.1 family inclusion and
two A.3.2 survivals prove
$W_p^{\rm rad}\le P_{p,2}\le P_p$.

To prove the upper bound for
\emph{every} parent leaf rather
than only circular ones, use the
Joukowski uniformization of the
plane complement of $\T_p$.
Set $m=(x_{2p}+x_p)/2$,
$r=(x_p-x_{2p})/2$ and
\[
 J(\zeta)=m+\frac r2(\zeta+\zeta^{-1}),
 \quad \zeta=-e^{s+it},\qquad
 a=\arccos\frac{m-x_{5p}}r,\quad
 b=\arccos\frac{m-x_{3p}}r,\quad0<a<b<\pi.
\]
The left-tip upper/lower
end-interval lifts in the
logarithmic cover are
$E_-=\{s=0,-b<t<-a\}$
and $E_+=\{s=0,a<t<b\}$.
Let $\widetilde\Gamma_p$
be \emph{all} proper
locally rectifiable curves
in the lifted physical
target joining these
sets.  KKL A.1
subordination to the
specified \emph{left-tip}
base selects precisely this
lift: a route around the
right tip, an additional
whole-slit winding, or
an outer-boundary endpoint
does not belong to that
relative-endpoint class.
Every admissible parent
leaf therefore lifts into
$\widetilde\Gamma_p$.
No circularity, simplicity
or child property is assumed
of $\widetilde\Gamma_p$.

One cannot equate conformal
width upstairs and downstairs
under a noninjective cover.
Here is the required
inequality proof.
For an admissible finite-area
density $\widetilde\rho$
on the lifted surface, define
on the physical surface
\begin{equation}\label{eq:period}
 \rho(\pi z)^2=
 \sum_{n\in\mathbb Z}
 \widetilde\rho(z+2\pi i n)^2.
\end{equation}
The sum is well-defined
almost everywhere; Tonelli
gives equal total areas,
and pointwise $\rho$ is at
least the chosen lifted
density along any selected
lift of a parent leaf.
Thus the pushdown density
is admissible for every
parent leaf.  Taking
infima proves
$P_p\le\W_{\widetilde\Om_p}
(\widetilde\Gamma_p)=:Q_p$.
This is an inequality, not
a width identity under
a covering.

The lifted physical
domain lies in the
upper full strip envelope
$0<s<L_+$, where
\[
 B_+=\frac{R_p+m}{r},\qquad
 \rho_+=B_++\sqrt{B_+^2+1},\qquad
 L_+=\log\rho_+.
\]
Indeed $|J(\zeta)|\ge
(r/2)(|\zeta|-|\zeta|^{-1})-m
>R_p$ when $|\zeta|>\rho_+$.
Conformal monotonicity
for the \emph{full lifted
path class}, followed by
exponentiating the strip
and applying the
Schwarz--Christoffel
quadrilateral formula
\cite[Chapter I, \S\S1--4]{AhlforsInvariants},
gives
\begin{align}
 Q_p&\le Q(L_+),&
 Q(L)&=\frac{K(k_L)}{K(\sqrt{1-k_L^2})},&
 k_L&=\frac{\sinh(\pi(b-a)/(2L))}
             {\sinh(\pi(b+a)/(2L))},\label{eq:Q}\\
 Q(L_+)&\le Q_{\rm plane,p},&
 Q_{\rm plane,p}&=\frac{K(k_\infty)}
                       {K(\sqrt{1-k_\infty^2})},&
 k_\infty&=\frac{b-a}{b+a}.\label{eq:Qplane}
\end{align}
Here $K(k)=\int_0^{\pi/2}
(1-k^2\sin^2\theta)^{-1/2}\,d\theta$;
the last inequality uses
$\sinh u/\sinh v\le u/v$
for $0<u<v$ and
monotonicity of $K(k)/K(k')$.
Equations \eqref{eq:Q}--\eqref{eq:Qplane}
are \emph{full all-path
upper envelopes}, not
claimed optimizer widths.

For $3\le p\le12$
the strict rational certificates
give $W_p^{\rm rad}>0.25$
and $Q(L_+)<0.78$,
with the individual finite
rows in the source audit.
For $13\le p\le16$
independent adjacent-dyadic
root cells give
$W_p^{\rm rad}>0.2604925$
and $Q(L_+)<0.771$.
For all $p\ge17$
the normalized-coordinate
error estimates yield
\begin{equation}\label{eq:tail-sandwich}
 W_p^{\rm rad}>
 0.2604925-7\,000\,000\,4^{-p}>0.26,\qquad
 Q_{\rm plane,p}<0.7750748<0.78.
\end{equation}
The first inequality is
indeed $>0.26$ at $p=17$
and improves thereafter.
For the upper error box,
$|\tau_p-\tau_*|\le15246q$
and the $j=2,3,5$ cut
errors are bounded by
$(30\cdot15246+88(4^j-1)/3)q$;
the box at $p=17$
contains every later one.
The exact rational-interval
cosine, logarithm and
elliptic AGM bounds
certify the numerical
claims.  As $p\to\infty$,
\eqref{eq:tail-sandwich}
gives a lower liminf at
least $0.2604925$ and
an upper limsup at most
$0.7750748$; both are
strictly inside $(0.26,0.78)$.
No continuity theorem
for the unrestricted
supremum has been used.
\end{proof}

\begin{theorem}[Strict gain of \emph{independent}
unbuffered Wave suprema]\label{thm:gain}
For every integer $p\ge3$,
$P_p,C_p^+,C_p^-$ are
positive and finite and
\begin{equation}\label{eq:gain}
 D_p:=\frac1{C_p^+}+\frac1{C_p^-}
             -\frac1{P_p}
   =\frac2{C_p}-\frac1{P_p}>\frac1{100}.
\end{equation}
The strict arithmetic cases
are
\begin{align}
 P_3&>49/100,& C_3&<487/500,&
 D_3&>300/23863>1/100,\nonumber\\
 P_4&>511/1000,& C_4&<197/200,&
 D_4&>7400/100667>7/100,\label{eq:gain-cases}\\
 P_p&>581/1000,& C_p&<247/250,&
 D_p&>43500/143507>3/10\quad (5\le p\le12),
 \nonumber\\
 P_p&>13/20,& C_p&<63/50,&
 D_p&>40/819>1/100\quad (p\ge13).\nonumber
\end{align}
These are strict inequalities for the
\emph{full} marked A.1 suprema;
they do not assume the suprema
are attained or paired.
\end{theorem}
\begin{proof}
Use the positive conformal
normalization $\xi=-A_pz$,
whose scale is $\lambda_p=-A_p>0$
(and distinguish it from
the negative-order normalization
used in the real recurrence).
Write $X_j=\lambda_px_{jp}$,
$\rho_p=\lambda_pR_p$,
$m=(X_2+X_1)/2$ and
$r=(X_1-X_2)/2$.  In the
target Joukowski logarithmic
coordinates the conformal map
to the physical target is
\[
 \mathcal{J}_P(-e^{s+it})=
 m-r\cosh(s+it).
\]
It is injective for $s>0$,
$|t|<\pi$; it meets the
left-tip electrodes at the
same two angles $0<a<b<\pi$
as in the sandwich proof.
We construct an \emph{embedded
full quadrilateral foliation}
of A.1 Waves for a lower
bound, not merely evaluate an
abstract all-path modulus.

For $p=3,4$ take the right
circular segment
$Q_{h,T}=\{\Re z>0,\ |z+d|<L\}$,
with $T=31/10$, with
$h=83/100$ or $88/100$
respectively, and
\[
 d=\frac{T^2-h^2}{2h},\qquad
 L=\frac{T^2+h^2}{2h}.
\]
The ratio $(z-iT)/(z+iT)$
maps it to a wedge of opening
$2\arctan(h/T)$; raising to
$\eta=\pi/(2\arctan(h/T))$
maps to a half-plane.
On the imaginary diameter
the two electrode values
are expressed by
$u=\eta\operatorname{artanh}(a/T)$,
$v=\eta\operatorname{artanh}(b/T)$.
The symmetric cross-ratio
modulus is
\begin{equation}\label{eq:segmod}
 k=\frac{\sinh(v-u)}{\sinh(v+u)},\qquad
 W_{\rm seg}=\frac{K(k)}{K(\sqrt{1-k^2})}.
\end{equation}
The full-domain inequalities
$b<T<\pi$ and
$|m|+r\cosh h<\rho_p$
are certified over the
selected \emph{entire}
complex quadrilateral,
so its SC rectangle
\emph{foliation} maps
to a family of simple
disjoint left-tip A.1
paths with base
$(I_9,I_0,I_1)$.
The exact interval
AGM evaluation
of \eqref{eq:segmod}
proves $P_3>49/100$
and $P_4>511/1000$.
A right semidisk of radius
$5/2$ would \emph{not}
satisfy the $p=3,4$
physical containment
and is not used.

For $5\le p\le12$
take the right semidisk
$\{\Re z>0,\ |z|<R\}$
with $R=7/5$.
The Möbius map
$z\mapsto2Rz/(R^2-z^2)$
takes it to a half-plane;
on the imaginary diameter
the real symmetric
electrodes have values
$B_R(t)=2Rt/(R^2+t^2)$.
The quadrilateral ratio
and embedded Wave width are
\begin{equation}\label{eq:semidiskmod}
 k=\frac{B_R(b)-B_R(a)}{B_R(b)+B_R(a)},
 \qquad W_{\rm semi}=\frac{K(k)}
                              {K(\sqrt{1-k^2})}.
\end{equation}
The \emph{whole} semidisk
is inside the target
because $b<R<\pi$
and $|m|+r\cosh R<\rho_p$
on each exact finite chart.
The SC rectangle foliation
therefore supplies a lower
\emph{actual A.1 parent
competitor} of width
$>581/1000$ for every
$5\le p\le12$.

For all $p\ge13$ use the
same construction with
$R=5/2<\pi$.
The largest normalized
orbit interval box is
at $p=13$; it proves
$b<R$, $|m|<2$, $r<2$,
$e^{5/2}<13$, and
\[
 |\mathcal{J}_P(-e^{s+it})|
 \le|m|+r\cosh R<16
 <\frac3{20}\,2^p
 <\lambda_pR_p.
\]
The last strict gate uses
\eqref{eq:A-bds}.
All later normalized boxes
are nested inside the $p=13$
box.  Exact evaluation of
\eqref{eq:semidiskmod}
therefore proves the
\emph{all-period} strict
embedded width $P_p>13/20$.
\emph{Integration note, 24 November 2026.}
This is the corrected single-box bound. The separate
\href{https://github.com/dovvnloading/loops-research-archive/blob/main/research/independent_audits/tip_family_wave_gain_correction_20261123/NOTE.md}{versioned Wave-gain correction},
Section~4, isolates the actual $p=13$ parameter and starts its uniform
box at $p=14$. It proves the stronger $P_p>2/3$, hence
$D_p>11/126$, for every $p\ge13$ with the same full error coefficient.
Both certificates coexist; the wider $p=13$ box used here does not
establish that sharper floor. Their physical interpretation still
requires the stated analytic class and metric arguments.
The embedding, injectivity,
correct homotopy and
physical gate are checked
independently of the
abstract SC formula.

For the \emph{upper} child
bound, let
\[
 a_c=\arccos(x_{4p}/x_p),\
 b_c=\arccos(-x_{2p}/x_p),\
 c_c=\arccos(-x_{5p}/x_p),\
 d_c=\arccos(\alpha_p/x_p).
\]
The strict cuts give
$0<a_c<b_c<c_c<d_c<\pi$.
Lift the entire source
complement of its slit
into the universal
right-half-plane chart
$\xi=x_p\cosh(s+it)$.
The \emph{full} lifted
four-electrode quadrilateral
has cross-ratio
\begin{equation}\label{eq:child}
 k_c^2=
 \frac{(b_c-a_c)(d_c-c_c)}
 {(c_c-a_c)(d_c-b_c)},\qquad
 E_p=\frac{K(k_c)}{K(\sqrt{1-k_c^2})}.
\end{equation}
By the square-sum metric
periodization \eqref{eq:period},
any admissible physical
child Wave width is at most
this \emph{all-path}
envelope.  In particular
no child maximal Wave is
identified with an
elliptic optimizer.

For $3\le p\le12$
there is a \emph{strict}
source-domain saving.
The actual source lies in
$|z|<R_p$.  In its log
Joukowski chart the open
unit box
\[
 B_p=(s_p,s_p+1)\times(0,1),\quad
 (s_p)_{p=3}^{12}=
 (2,2,2,3,4,4,5,6,6,7)
\]
lies \emph{outside} this
source because the
rational Taylor lower
bound certifies
$x_p\sinh(s_p)>R_p$,
while
$|\cosh(s+it)|^2=
 \sinh^2s+\cos^2t$.
An admissible extremal
metric for the full
half-plane rectangle
can be periodized and
then restricted to the
actual source.  It remains
admissible for every
physical child path,
but its area is smaller
by the positive energy
of this entire omitted
box.  The four-electrode
SC derivative, bounded
below on the box, and
the rectangle height
bound certify
\[
 \int_{B_p}\rho_{\rm SC}^2\,dA>
 \frac1{5\,000\,000},\qquad
 C_p\le E_p-\frac1{5\,000\,000}.
\]
This is an \emph{actual
physical source metric}
saving, not the fictitious
claim that a target-parent
semidisk fits inside
the source.  The
$2^{-512}$ outward
evaluations of
\eqref{eq:child} and
the missing energy give,
respectively,
$C_3<487/500$,
$C_4<197/200$ and
$C_p<247/250$ for
$5\le p\le12$.

For $p\ge13$ a simpler,
weaker estimate suffices:
the uniform normalized
orbit/cut enclosure
around the model root
$\tau_*$, with
$|\tau_p-\tau_*|\le15246q$,
contains every actual
child cross-ratio in
the $p=13$ worst-case
interval box.
The complete full
right-half-plane
elliptic envelope
\eqref{eq:child} is
strictly $<63/50$
there, using rational
cosine inverses and
AGM remainder bounds.
Consequently $C_p<63/50$
\emph{for every} $p\ge13$,
without any missing-box
assumption at large $p$.

Finally take reciprocals
in the \emph{correct}
direction.  From $P_p>L_p$
and $0<C_p<U_p$,
\[
 \frac2{C_p}-\frac1{P_p}>
 \frac2{U_p}-\frac1{L_p}.
\]
Substitution yields
\[
 \frac2{487/500}-\frac1{49/100}
  =\frac{300}{23863},\quad
 \frac2{197/200}-\frac1{511/1000}
  =\frac{7400}{100667},
\]
\[
 \frac2{247/250}-\frac1{581/1000}
  =\frac{43500}{143507},\quad
 \frac2{63/50}-\frac1{13/20}=\frac{40}{819}.
\]
Each is strictly greater than the
claimed threshold.  All three
suprema are defined on their
independent full classes;
no common parent Wave has
been selected in this argument.
\end{proof}

\begin{proposition}[Leafwise-pairable loss, not a contradiction]
\label{prop:loss}
Let $P_{p,\pp}$ be the
supremum over \emph{only}
those parent A.1 Waves
whose every leaf has both
named A.3.2-surviving
proper source restrictions,
for which $\iota$ is injective
on their union, and for
which the two restrictions
belong to the two
designated source classes.
Then
\[
 P_{p,\pp}\le C_p/2,\qquad
 \frac2{C_p}-\frac1{P_{p,\pp}}\le0.
\]
\end{proposition}
\begin{proof}
For each such parent
$\mathcal S$, KKL
\cite[Lemma A.11 and
Proposition A.12]{KKL2026}
gives, with its
\emph{properness and
injectivity hypotheses},
\[
 \W(\mathcal S)\le
 \W(\mathcal S_+)\oplus\W(\mathcal S_-),
 \qquad a\oplus b=\frac{ab}{a+b}.
\]
By definition each
source width is at most
the independently optimized
$C_p$, so their harmonic
combination is $\le C_p/2$.
Taking the supremum over
the \emph{smaller}
pairable class gives
the first inequality,
and reciprocal inversion
gives the second.
In particular
Theorem~\ref{thm:gain}
proves $P_p>P_{p,\pp}$.
It \emph{does not}
provide gain for
the two restrictions
of any one parent.
\end{proof}
